\section{Protection of vulnerable people: a duty of care, written down}
\label{sec:protection}

The sixth appeal is protective rather than purely compassionate. Health debates
return again and again to the groups most exposed to harm: children, the elderly,
disabled patients, veterans, and rural communities. For medical AI the specific
worries are bias against minority populations, errors that fall hardest on the
vulnerable, and consent that is assumed rather than confirmed. The same bias that
the previous section measured falls most heavily on the people least able to
absorb it \cite{obermeyer2019dissecting}, which is why a national action plan for
AI as a medical device places real-world performance monitoring across subgroups
at its center \cite{fda2021samd}.

Protection is strongest when it is built into the procedure rather than promised
afterward. The author's site-documentation and consent work, including an adaption
of 21 CFR Part 50 for end-to-end Physical AI trials, makes informed consent a
verified, logged step rather than a signature gathered in haste
\cite{Kawchak2026PhysicalAIOncologyClinical, Kawchak2026Adaption21CFRPartb}.
Figure 8 shows the three gates a subgroup passes before its
results are accepted.

\begin{narrativefig}
\node[nfone] (a) at (0,0) {Subgroup data};
\node[nfdec] (b) at (3.4,0) {Bias parity};
\node[nfdec] (c) at (6.8,0) {Consent verified};
\node[nfdec] (d) at (10.2,0) {Outcome audit};
\node[nfact] (e) at (8.5,-2.1) {Accepted with report};
\node[nfrisk] (f) at (4.0,-2.1) {Held for human review};
\draw[nfedge] (a)--(b); \draw[nfedge] (b)--(c); \draw[nfedge] (c)--(d);
\draw[nfedge] (d) to[out=-90,in=0] (e.east);
\draw[nfedged] (b) to[out=-90,in=90] (f.north);
\end{narrativefig}
\vspace{-0.15cm}
\figcaption{Figure 8. Safeguards written into the procedure (Mermaid 15 as colored
TikZ): parity, consent, and audit each guard the accept, or hold the work for a
human.}

\autoref{tab:protection} lists each protected group, its specific risk, and the
safeguard that answers it.

\needspace{9\baselineskip}\begin{table}[H]
\footnotesize
\begin{tabularx}{\textwidth}{@{}L{3.0cm} Y L{4.0cm}@{}}
\toprule
\textbf{Group} & \textbf{Specific risk} & \textbf{Safeguard}\\
\midrule
Children & Too few cases for ordinary evidence & Subgroup audit and reporting\\
Elderly and disabled & Models trained on others & Bias-parity surveillance\\
Veterans & Fragmented records & Verified, auditable record\\
Rural communities & Distance and access gaps & Portable, standardized platform\\
\bottomrule
\end{tabularx}
\caption{Protected groups, their risks, and the bill safeguards.}
\label{tab:protection}
\end{table}

Protection that is written into statute, measured, and reported is stronger than
protection that depends on the good intentions of whoever happens to be on duty.
The quiet question is gentle but firm. If we can guarantee the check in law, why
would we leave the most vulnerable patients relying on hope alone?
